\section{Conclusions}
\label{sec:conclusiones}
\begingroup
\fontsize{10.5}{13}\selectfont
\setlength{\parskip}{2pt plus 1pt minus .5pt}

La construction détermine conjointement \(\pi,\varphi,e\) à une profondeur
arbitraire et dérive \(\alpha\) comme coordonnée de clôture des régions
et du registre dodécaphasique \(K\). La même structure produit la
composition électronique, sa représentation de spin \(1/2\) et les
projecteurs d'hélicité, et conserve l'incidence qui conduit à Leech
et à \(V^\natural\), avec le Monstre pour groupe d'automorphismes.
Ces résultats procèdent des évaluations additive et multiplicative
d'APP avec leurs résidus et leurs quotients, de l'orientation et des régimes
quadratiques de TRIT, ainsi que de la sélection, du transport et de la mise à jour
de la mémoire du TPK. Leurs domaines, leurs normalisations et leurs applications de
compatibilité déterminent tant les valeurs que les relations qui
persistent entre leurs réalisations.

L'exhaustivité initiale et la profondeur du raffinement sont des propriétés complémentaires. Les \(104\,976\) conditions orientées déterminent \(468\) émissions, dont la réduction contient \(243\) mots visibles dans l'espace ambiant ternaire à \(729\) éléments. Sur ce domaine se distinguent les régions de clôture, de propagation et d'auto-échelle. La région de clôture conserve cinq émissions de multiplicités \(432+4\cdot144=1008\). Leurs réalisations archimédiennes se déterminent au moyen de cylindres compatibles de diamètre décroissant. Tout préfixe fini s'obtient par un calcul fini, tandis que la compatibilité identifie l'objet complet. Le retour nonadique conserve la phase et augmente la mémoire.

La réalisation sturmienne précise l'ordre apériodique de ce transport. Les séparations primaires \(21/22\), les retours secondaires \(6/7\) et les complexités \(m+1\), \(9(m+1)\) et \(3(m+1)\) proviennent d'une même pente et de ses lectures de phase. Les relations de Weyl et l'incrément de mémoire articulent la rotation de base avec l'algèbre des observables et des translations.

La formulation au moyen de résolvantes explicite l'accumulation de mémoire : la série génératrice du registre affine s'obtient en composant la résolvante du transport avec la série des incréments. La composition de trois segments conserve exactement le cocycle. Lorsqu'un sous-espace auxiliaire est éliminé, le complément de Schur détermine la correction de l'opérateur effectif et, simultanément, la modification de la source et de la fonctionnelle d'évaluation. Ces identités sont démontrées dans le corps de l'article et conservent la contribution des états qui interviennent entre deux observations.

Le registre dodécaphasique introduit une autre forme de conservation. Les canaux de la trajectoire produisent le vecteur entier signé \(U\), et la transformation par orbites détermine
\[
 K=\tfrac14\mathcal H_{12}U,\qquad U=\mathcal H_{12}K.
\]
Ces équations expriment une équivalence réversible entre coordinatisations entières, au sein du sous-réseau spécifié. Les différences de pas trois et quatre, avec la somme totale, permettent de reconstruire le même vecteur. La réversibilité correspond à ces registres ; une lecture ultérieure des supports conserve une information différente.

La réalisation rationnelle périodique
\[
 \kappa_{\mathrm{per}}=\frac{N_K}{1000^{12}-1}
\]
évalue le registre canonique avec une origine et une orientation fixées. Sa période minimale de douze blocs explique le dénominateur et distingue cette constante de la lecture finie \(N_K/1000^{12}\). L'objet structurel \(K\), son équation de reconstruction et le nombre rationnel qui le représente restent distincts. La périodicité appartient à cette évaluation, tandis que les coordonnées régionales et leurs retenues maintiennent leurs propres conditions de prolongement.

La détermination positionnelle de \(\alpha\) compose les registres régionaux avec \(K\) au moyen de la normalisation positionnelle. Les blocs et les quotients transportés satisfont une égalité locale dont la somme télescopique conserve la frontière. La section complète détermine
\[
 \alpha_A=\pi+e-\varphi-4-\kappa_{\mathrm{per}}.
\]
Le quatre provient des parties entières régionales ; l'orientation détermine les signes ; le prolongement compatible fixe la retenue. L'équation exprime la coordonnée de compatibilité arithmétique de la clôture. Son évaluation commence par
\[
 \alpha_A=0.007297352569283800997285105472380662999\ldots.
\]

La représentation analytique corrélée conserve le jet d'ordre neuf
et spécifie son prolongement par
\(b_{10+3n}=\lambda/729^n\), \(b_{11+3n}=b_{12+3n}=0\).
Le coefficient \(\lambda\) est calculé à partir des coordonnées régionales et
terminales du même état, au sein de la classe géométrique déclarée.
La convergence uniforme de la fonction et de sa dérivée, l'unicité de la
racine et les intervalles rationnels compatibles avec les cylindres entiers
établissent l'égalité des préfixes des deux représentations à toute profondeur
(théorème~\ref{thm:alpha-dos-publicaciones-completas}). La proximité entre
la racine de la troncature et \(\alpha_A\) reste une comparaison finie ;
l'égalité complète découle de la récurrence et de sa preuve de
compatibilité, non de cette proximité ni d'une valeur métrologique.
La coordinatisation analytique est une représentation du même point engendré,
et non une seconde sélection causalement indépendante.

La construction électronique reprend le centre d'APP après les coordonnées qu'elle utilise. Le point fixe \((5,5)\) et les déformations \((8,2)\), \((2,8)\) conservent la somme et le résidu multiplicatif, tout en distinguant une unité de quotient. Leur représentation matricielle explique la norme \(\sqrt3/4\). La coordonnée basale compose cette norme avec le transfert régional, le couplage et la normalisation torsionnelle ; la mémoire de retour agit ensuite multiplicativement.

Deux lecteurs explicitent ce retour : l'un compte les discordances chronologiques et l'autre compare les corrélations de fenêtres dodécaphasiques. Leurs domaines et leurs opérations sont différents, mais sur la trajectoire électronique, tous deux déterminent \((80,54,6)\). La composition est
\[
 \varepsilon_e^\beta
 =\frac{\sqrt3}{4}
 \left[\exp\!\left(\frac{\varphi}{\pi^2}\right)-22\alpha^3\right]
 (1+15\Delta_4)
 R_{\mathrm{act}}^{\,80-54\alpha+6\Delta_4},
\]
dont l'évaluation commence par \(0.510998950690000808608\ldots\). L'égalité centrale des triplets relie ces représentations sans identifier universellement leurs lecteurs. L'échelle commune d'action se simplifie dans \(R_{\mathrm{act}}\) ; la réalisation en énergie ou en masse ajoute la coordinatisation dimensionnelle. Le quantum chronologique conserve sa fonction d'échelle et la section électronique conserve sa trajectoire propre.

Le transport central détermine en outre un support exact de trois diagonales
cycliques, comportant \(27\) cellules. Les deux orientations produisent des registres
d'événements ayant le même histogramme de valeurs absolues et les mêmes signes
sectoriels, mais des marges distinctes
dans la coordinatisation intégrale réversible. La lecture harmonique du registre central
concentre ses puissances dans les modes \(2,6,10\), de valeurs \(64,16,64\).
La pondération démontrée restitue \(\Omega=80\) et \(P_6=6\), et exprime
l'exposant de retour de l'équation électronique comme une fonctionnelle
spectrale. L'histogramme, l'autocorrélation et le registre réversible
conservent des degrés d'information différents sur la trajectoire.

La représentation spinorielle de la section~\ref{sec:el-representacion-espinorial} explicite l'information rotationnelle du bloc électronique. Les deux feuillets induisent un triplet hermitien et ses générateurs de rotation ; pour chaque direction, les projecteurs spectraux séparent deux droites complexes orthogonales. La loi de transport distingue le changement de signe après une rotation complète et le retour spinoriel après deux rotations. Lorsque la direction est réalisée comme direction de propagation, cette décomposition détermine l'hélicité. L'inversion de direction, la conjugaison de charge, la chiralité et l'holonomie nonadique conservent leurs domaines respectifs. L'évaluation électronique agit scalairement sur la représentation, de sorte que cette structure rotationnelle est conservée sans modifier la formule de masse ni ses lecteurs.

La prolongation exceptionnelle développe l'incidence du même état. Le code ternaire conserve les mots ; leurs supports retiennent les positions et les intersections. Le drapeau marqué de \(K\) et du registre antérieur à la retenue de \(\alpha\) relie les deux lectures. Le relèvement binaire relie les quatre pétales de Witt à une cinquième octade transversale. Le recollement ternaire produit le réseau pair unimodulaire, et le voisin sélectionné conduit à Leech. Avec les données algébriques déclarées, la construction de Frenkel--Lepowsky--Meurman donne \(V^\natural\), dont le groupe d'automorphismes est le Monstre\cite{flm}. Moonshine décrit ensuite ses traces graduées\cite{borcherds}. L'action correspond à l'algèbre obtenue ; la construction numérique et cette réalisation exceptionnelle restent reliées par leurs applications à partir de l'état commun.

La comparaison métrologique s'effectue à la fin\cite{codata2018,codata2022}. \(\alpha\) coïncide après arrondi avec le centre de CODATA 2018 et s'écarte d'environ \(4.53\) incertitudes-types du centre de 2022. L'expression électronique coïncide après arrondi avec le centre de 2022 et s'écarte d'environ \(4.60\) incertitudes-types de celui de 2018. L'exactitude algébrique, l'erreur d'évaluation et l'incertitude expérimentale ont des significations différentes. Ni la coïncidence avec une formule ni la date d'un ajustement ne décident à elles seules de sa proximité avec la valeur physique.

La comparaison historique supplémentaire du
paragraphe~\ref{sec:contraste-ventana-reciproco} utilise la fenêtre finie
\(\alpha_{12}^{(0)}\) et l'inverse du centre
\(R_{2018}=137.035999084\). La différence dans la coordonnée inverse est
approximativement \(1.066588006664435\times10^{-32}\), avec une différence
relative \(7.783268730800000\times10^{-35}\) ; dans la coordonnée directe,
elle est approximativement \(-5.679725607012965\times10^{-37}\).
Cette proximité correspond à la troncature décimale de l'inverse
imprimé. Sa précision arithmétique, la comparaison avec le centre tabulé
de \(\alpha\) et l'incertitude expérimentale conservent leurs domaines
respectifs.


La direction \(u_K=P_3P_{11}K\) étend cette description de la fonction de \(K\). Sa norme exacte positive démontre que le registre sélectionne une droite, et le projecteur orthogonal complémentaire réalise la réduction \(11\to10\). La réduction élimine une direction spécifiée par le registre ; l'application coordonnée conserve en outre l'écran réticulaire \(26\to24\) et entrelace l'échange de rayons réciproques. Chaque facteur maintient son propre domaine et l'information qu'il conserve.

La construction exceptionnelle admet également deux réalisations par
orbifolds, d'ordres \(2\) et \(3\), sur le même voisin de réseau.
Les identifications avec \(V^\natural\) fournissent un isomorphisme
qui conserve le vide, le vecteur conforme, le produit de vertex et
la graduation. Les traces de Fricke et l'automorphisme dual de l'orbifold
décrivent des opérations concrètes sur ces réalisations. L'exposant
\(-1/12\) du quotient de fonctions êta de niveau trois possède ici
une place précise, distincte d'une somme ordinaire divergente.
Ces applications relient les deux réalisations au moyen de morphismes
qui conservent leurs données algébriques et leurs identifications classiques.

% BEGIN S27_FRAMING
Le cocycle de retenue détermine comment se conserve l'évaluation entière
lors de la composition des colonnes régionales. Les identités
\(p+e+f=q+3c\) et \(p+e-f=a+3h\) retrouvent les deux combinaisons
entières à partir de leurs résidus et quotients. Le tableau complet montre
que \((0,0,0)\) et \((1,2,0)\) partagent \(q=a=0\), tandis que leurs
quotients sont respectivement \((c,h)=(0,0)\) et \((1,1)\).
Les représentations polynomiales et les rangs \(4,5,5\) précisent
la contribution algébrique de ces retenues. Les neuf fibres de trois
triplets de l'application résiduelle \((p,e,f)\mapsto(q,a)\) conservent explicitement
la liberté de répartition entre les entrées, et les récurrences
transportent les évaluations conjointes et signées dans le domaine
entier déclaré.

La lecture d'incidence distingue les trois étoiles régionales par les signatures \((3,1,0)\), \((1,0,3)\) et \((1,1,2)\), avec leurs douze complétions explicites. Le recensement global répartit les \(495\) quadruplets en six classes, et la restriction régionale classe les \(54\) hexades pertinentes. La succession de filtres sur le groupe construit prouve que le repère ordonné \((H^-_\alpha,H_e,H_\varphi,H_{\rm APP})\) a un stabilisateur trivial. Cette rigidité détermine la liberté résiduelle de son action sur les douze positions. La récupération numérique de \(K\) conserve ses propres équations réversibles et le registre entier complet. Deux fonctions de la même généalogie se trouvent ainsi articulées : grandeurs, résidus et quotients dans l'évaluation arithmétique ; supports, polarisation et stabilisateurs dans la réalisation d'incidence.

% END S27_FRAMING

% BEGIN R27_FRAMING_CONCLUSION_ES
Les preuves de l'émetteur et de sa prolongation précisent les degrés
d'information que conserve cette généalogie. Les deux signatures déterminent les
\(42\) blocs décimaux possibles d'une fenêtre ; six émissions témoins
établissent le cycle des raccordements régionaux. Un histogramme commun peut
s'accompagner de corrélations temporelles distinctes, de sorte que l'archive
ordonnée conserve une distinction opératoire supplémentaire. Les contrôles de
jauge et de frontière identifient, respectivement, la normalisation de
l'observable et la dépendance à l'égard du sens de retenue, de l'extrémité
terminale et de la position des coefficients.

Dans APP, les lois du défaut maintiennent explicites le module, le support et la
dimension. La restriction de l'opérateur agrégé à un sous-réseau
d'indice deux produit \((-1)\oplus3H\), avec \(\det H=1\). La récurrence
de Pell, avec les données initiales fixées, détermine les puissances de
\(H\), qui conservent une norme égale à un. Le codage nonadique conserve
les six trits au moyen de trois paires et transporte la somme avec retenue.
La corrélation quadratique donne \(C^2=5I_6\) ; la saturation, la neutralité
duale et l'orientation fixent par étapes la frontière de prolongement.
Enfin, l'identité de niveau trois démontre
\(J-(t_3+12)\in3^9q\mathbb Z[[q]]\), avec un exposant uniforme optimal.
Ces résultats relient les opérations finies, leurs lois générales
démontrées et la réalisation exceptionnelle, tout en maintenant le domaine propre
de chaque affirmation.
% END R27_FRAMING_CONCLUSION_ES

% BEGIN R28_FRAMING_CONCLUSION_ES
Les cinq clôtures régionales atteignent le minimum huit dans le problème fini de couverture des quatre arêtes prescrites. Le graphe conserve les étiquettes des émissions et le quotient conserve leurs incidences entre classes de phase. L'inversion orientée rend impair le lecteur \(\nu_{120}\) et maintient pair \(\nu_{270}\) ; ces identités explicitent la réponse de la mémoire au changement d'orientation.
% END R28_FRAMING_CONCLUSION_ES
\par % S27_LAYOUT_CONCLUSION; REV03: continuous conclusion pagination.
La complétion cubique et le transport apériodique précisent deux sens
de la conservation qui s'articulent dans le développement. La première conserve
l'unité par projection de mesures normalisées et distingue le cube
initial, la complétion perforée de \(999\) éléments et la restitution
du sommet. Le second conserve le registre de trajectoire durant le
retour de phase. La monodromie et le modèle mathématique de cristal temporel
apériodique expriment ce second mécanisme au moyen d'applications et
d'observables, et non seulement par une image hélicoïdale. L'exponentielle
tritique, quant à elle, unifie les réalisations elliptique, parabolique et
hyperbolique sans supprimer leurs relations quadratiques respectives.

La constante de Planck réduite, \(\hbar\), occupe une position explicite dans cette architecture. Le quotient déterminantal d'ordre \(16\) agit dans la norme qui détermine la décade \(-34\) ; le caractère \(H_5\) et sa continuation régionale déterminent les deux sections d'action dont le quotient entre dans la coordonnée électronique. La comparaison de la section précédente avec \(h/(2\pi)\) montre une différence relative de l'ordre de \(10^{-15}\), et la section retournée conserve la compensation qu'utilise l'électron. L'échelle d'action relie ainsi la norme déterminantale et le caractère régional au quotient électronique : la compensation appartient à la construction exacte, tandis que la proximité de la section précédente avec \(h/(2\pi)\) est une comparaison numérique et non une égalité exacte.


La conséquence conceptuelle se concentre dans l'articulation de \(K\),
de \(\alpha\) et de la section électronique. La reconstruction réversible de
\(K\) conserve des données intervenant tant dans la normalisation positionnelle
que dans l'incidence de frontière. La composition électronique utilise les
coordonnées précédemment construites et maintient son centre, ses transports
et sa réalisation dimensionnelle. Le prolongement exceptionnel conserve d'autres
relations du même état dans des représentations successives. Les connexions
entre ces objets font donc partie du résultat, au même titre que les
valeurs de leurs coordonnées.

Cette organisation détermine aussi l'ordre de la confrontation. Lorsqu'une
dérivation fixe une valeur sans la recevoir comme argument, la mesure
met à l'épreuve la conséquence obtenue. La correspondance avec une
observable physique exige, en outre, d'identifier la grandeur, la normalisation
et la réalisation dimensionnelle employées. En ce sens précis, la thèse
examine le passage d'une description paramétrée à l'explication des
relations qui déterminent ses paramètres. La portée de chaque résultat
est celle des démonstrations exposées : structure, équation et valeur
constituent des aspects inséparables de la détermination confrontée à l'expérience.

L'unicité du registre conserve deux recensements prospectifs et leurs
opérations propres (paragraphe~\ref{act21:sec:relacion-censos}).
La chaîne \(40320\to21902\to19446\to79\to1\) aboutit au
survivant hétérotypique de la face exceptionnelle ; la chaîne
\(196\cdot197\cdot196\to331\to2\to1\) aboutit à la sélection
orientée à partir des catalogues régionaux. Les populations intermédiaires
conservent le sens de chaque filtre. La composition des sorties
sélectionnées avec l'inversion entière articule détermination
prospective et récupération réversible.


En concluant sa conférence Nobel de 1946, Pauli a relié l'explication de l'électricité à une théorie capable de \emph{“determine the value of the fine-structure constant”}\cite{pauli1946} : déterminer la valeur de la constante de structure fine. La réponse mathématique obtenue ici est l'identité de clôture de \(\alpha\), sa représentation analytique compatible et sa composition avec la structure électronique. La génération régionale, la reconstruction réversible de \(K\) et les transports d'incidence font des valeurs et de leurs relations les conséquences d'une même construction. La métrologie compare ensuite leurs réalisations physiques ; elle ne fixe pas rétrospectivement les états ni les coefficients qui les déterminent.



\par\endgroup
